\clearpage\pagestyle{otMain}\setlength{\otBlocEspacement}{2.5mm}\fontsize{11}{15}\selectfont\corrigesChapitre{13}{Angles et droites parallèles}
\begin{blocLivre}{cor-13-1}\begin{corrige}{1}{Nommer un angle}\demarche{Lire la notation et la mesure.}
Le sommet est la lettre du milieu, $\res O$. Comme $0<40<90$, l’angle est $\res{\text{aigu}}$.
\label{fig-cor-13-1}\figureGeo[corrige]{ang-1}
\end{corrige}\end{blocLivre}
\begin{blocLivre}{cor-13-2}\begin{corrige}{2}{Un angle droit partagé}\demarche{Utiliser la mesure de l’angle total.}
Le total vaut $90^\circ$, donc $x=90-35=\res{55^\circ}$.
\label{fig-cor-13-2}\figureGeo[corrige]{ang-2}
\end{corrige}\end{blocLivre}
\begin{blocLivre}{cor-13-3}\begin{corrige}{3}{Un angle plat partagé}\demarche{Compléter jusqu’à 180 degrés.}
Ils sont supplémentaires : $180-128=\res{52^\circ}$.
\label{fig-cor-13-3}\figureGeo[corrige]{ang-3}
\end{corrige}\end{blocLivre}
\begin{blocLivre}{cor-13-4}\begin{corrige}{4}{Opposés par le sommet}\demarche{Reconnaître la paire.}
Les angles opposés par le sommet sont égaux : $\res{63^\circ}$.
\label{fig-cor-13-4}\figureGeo[corrige]{ang-4}
\end{corrige}\end{blocLivre}
\begin{blocLivre}{cor-13-5}\begin{corrige}{5}{Bissectrice}\demarche{Partager l’angle en deux parts égales.}
Les deux mesures sont égales et totalisent $74^\circ$ : $74:2=\res{37^\circ}$.
\label{fig-cor-13-5}\figureGeo[corrige]{ang-5}
\end{corrige}\end{blocLivre}
\begin{blocLivre}{cor-13-6}\begin{corrige}{6}{Quatre angles autour d’un point}\demarche{Combiner opposition et supplémentarité.}
L’opposé vaut $47^\circ$. Les deux adjacents valent $180-47=133^\circ$. Les quatre mesures alternent $\res{47,133,47,133}$.
\label{fig-cor-13-6}\figureGeo[corrige]{ang-6}
\end{corrige}\end{blocLivre}
\begin{blocLivre}{cor-13-7}\begin{corrige}{7}{Alternes internes}\demarche{Vérifier la position de la paire.}
Les droites sont parallèles, donc les angles alternes internes sont égaux : $\res{a=68^\circ}$.
\label{fig-cor-13-7}\figureGeo[corrige]{ang-7}
\end{corrige}\end{blocLivre}
\begin{blocLivre}{cor-13-8}\begin{corrige}{8}{Correspondants}\demarche{Identifier des positions correspondantes.}
Des angles correspondants découpés par une sécante sur deux parallèles sont égaux : $\res{112^\circ}$.
\label{fig-cor-13-8}\figureGeo[corrige]{ang-8}
\end{corrige}\end{blocLivre}
\begin{blocLivre}{cor-13-9}\begin{corrige}{9}{Prouver un parallélisme}\demarche{Partir de l’égalité des angles.}
Les angles alternes internes sont égaux. Par la propriété réciproque, les deux droites sont $\res{\text{parallèles}}$.
\label{fig-cor-13-9}\figureGeo[corrige]{ang-9}
\end{corrige}\end{blocLivre}
\begin{blocLivre}{cor-13-10}\begin{corrige}{10}{Exclure un parallélisme}\demarche{Utiliser une condition nécessaire.}
Non : si elles étaient parallèles, ces angles seraient égaux. Comme $\res{75\ne78}$, elles ne sont pas parallèles.
\label{fig-cor-13-10}\figureGeo[corrige]{ang-10}
\end{corrige}\end{blocLivre}
\begin{blocLivre}{cor-13-11}\begin{corrige}{11}{Deux étapes}\demarche{Compléter puis transporter un angle.}
Au premier croisement, l’angle adjacent aigu vaut $180-124=56^\circ$. Son correspondant au second croisement vaut donc $\res{56^\circ}$.
\label{fig-cor-13-11}\figureGeo[corrige]{ang-11}
\end{corrige}\end{blocLivre}
\begin{blocLivre}{cor-13-12}\begin{corrige}{12}{Équerres}\demarche{Additionner des angles adjacents.}
Le total vaut $30+45=\res{75^\circ}$. Son supplément vaut $180-75=\res{105^\circ}$.
\label{fig-cor-13-12}\figureGeo[corrige]{ang-12}
\end{corrige}\end{blocLivre}
\begin{blocLivre}{cor-13-13}\begin{corrige}{13}{Un angle à reconstituer}\demarche{Utiliser le tour complet.}
Un tour complet vaut $360^\circ$. Donc $x=360-110-95-80=\res{75^\circ}$. Contrôle : les quatre angles totalisent 360 degrés.
\label{fig-cor-13-13}\figureGeo[corrige]{ang-13}
\end{corrige}\end{blocLivre}
\begin{blocLivre}{cor-13-14}\begin{corrige}{14}{Justifier sans mesurer}\demarche{Citer la donnée et la propriété.}
Une impression ou une mesure ne prouve pas. Les deux angles correspondants codés sont droits, donc égaux à $90^\circ$. Par la propriété réciproque, $\res{(d)\parallel(e)}$ : deux droites perpendiculaires à une même droite sont parallèles.
\label{fig-cor-13-14}\figureGeo[corrige]{ang-14}
\end{corrige}\end{blocLivre}
\begin{blocLivre}{cor-13-15}\begin{corrige}{15}{Une bissectrice dans des parallèles}\demarche{Enchaîner supplément, correspondance et partage.}
L’angle aigu vaut $180-106=74^\circ$. Le parallélisme conserve cette mesure au second croisement. La bissectrice donne $74:2=\res{37^\circ}$ par moitié.
\label{fig-cor-13-15}\figureGeo[corrige]{ang-15}
\end{corrige}\end{blocLivre}
